(1) 对每支 \(\sigma=\pm1\) 定义平衡速度 \(u_\sigma=\sigma v_0\)、平衡密度 \(n_b=n_0/2\)；电子电荷 \(-e\)，其中 \(e>0\)。线性连续性与动量方程为
\[
(\partial_t+u_\sigma\partial_x)\delta n_\sigma+n_b\partial_x\delta v_\sigma=0,\qquad
(\partial_t+u_\sigma\partial_x)\delta v_\sigma=-\frac{e}{m}E .
\]
代入约定的时间空间谐波，得到
\[
\widehat{\delta v_\sigma}=-\frac{ie\widehat E}{m(\omega-ku_\sigma)},\qquad
\widehat{\delta n_\sigma}=-\frac{ike n_b\widehat E}{m(\omega-ku_\sigma)^2}.
\]
Poisson 方程 \(ik\widehat E=-e(\widehat{\delta n_+}+\widehat{\delta n_-})/\epsilon_0\) 在 \(\widehat E\ne0\) 下给出
\[
1-\frac{\omega_p^2}{2}\left[\frac1{(\omega-a)^2}+\frac1{(\omega+a)^2}\right]=0.
\]
该式要求 \(\omega\ne\pm a\)，对应无压冷流的共振奇点；不可把乘分母后的多项式所有零点一律认作本征模。

(2) 合并分式得 \((\omega^2-a^2)^2-\omega_p^2(\omega^2+a^2)=0\)。即
\[
X^2-(2a^2+\omega_p^2)X+a^2(a^2-\omega_p^2)=0,\quad
X_\pm=a^2+\frac{\omega_p^2}{2}\pm
\frac{\omega_p}{2}\sqrt{\omega_p^2+8a^2}.
\]
显然 \(X_+>0\)，且当且仅当 \(0<a<\omega_p\) 时 \(X_-<0\)，此时 \(\omega=\pm i\gamma\)，一支增长、一支衰减；当 \(a>\omega_p\) 两支均给稳定实频率。在 \(a=\omega_p\) 处 \(X_-=0\) 为稳定边界。特别 \(a=0\) 原式为 \(1-\omega_p^2/\omega^2=0\)，只有 \(\omega=\pm\omega_p\)；多项式额外引入的 \(X=0\) 是赝根，并不是物理静电波。

(3) 在不稳定区间，增长率为
\[
\gamma^2(a)=\frac{\omega_p}{2}\sqrt{\omega_p^2+8a^2}
-a^2-\frac{\omega_p^2}{2}.
\]
令 \(z=a^2\)，则 \(d\gamma^2/dz=2\omega_p/\sqrt{\omega_p^2+8z}-1\)，二阶导数
\(-8\omega_p/(\omega_p^2+8z)^{3/2}<0\)。因此唯一内部最大值在
\(z=3\omega_p^2/8\)，对应
\[
k_{\max}=\sqrt{\frac38}\frac{\omega_p}{v_0},\qquad
\gamma_{\max}=\frac{\omega_p}{2\sqrt2}.
\]
两端点增长率均为零，与极大值位置及稳定区间相符；只求使 \(\gamma\) 增大的群速而不验证根号分支将错误地漏掉这一不稳定性。

(4) 若 \(0<a\ll\omega_p\)，利用
\(\sqrt{1+8a^2/\omega_p^2}=1+4a^2/\omega_p^2-8a^4/\omega_p^4+\cdots\)，得到
\(\gamma=a[1-2a^2/\omega_p^2+O(a^4/\omega_p^4)]\)。
若 \(a\gg\omega_p\)，有
\[
\omega_{\pm,\mathrm{positive}}=
a\pm\frac{\omega_p}{\sqrt2}+O(\omega_p^2/a),
\]
对应随束流漂移的两支稳定波（两种正负频率另由时间反演给出）。令
\(W=\epsilon_0E^2/2+\sum_\sigma mn_b(\delta v_\sigma)^2/2\)，将线性动量方程分别乘以 \(mn_b\delta v_\sigma\)，并与纵向 Maxwell–Ampere 式
\(\epsilon_0\partial_tE=e\sum_\sigma(n_b\delta v_\sigma+u_\sigma\delta n_\sigma)\)
相加，精确获得
\[
\partial_tW+\partial_x\sum_\sigma\frac{mn_bu_\sigma}{2}(\delta v_\sigma)^2
=eE\sum_\sigma u_\sigma\delta n_\sigma.
\]
右侧是扰动与背景有向流交换能量的项；仅靠上式的正定平方项并非完整的非线性能量。若背景束流允许反作用，其二阶平均流能变化补偿该交换，全系统能量仍守恒。热速度、碰撞、有限束宽和动力学 Landau 共振均不包含在此冷且无界的理想化模型中。